\slajdid{02-s032}
\begin{frame}[label=02-s032]{Minimizacija logičkih funkcija}
% element: 02-s032:e05
U prethodnom delu razmatrali smo mogućnost realizacije kombinacionih mreža. Pojavom digitalnih kola od interesa je postala i \alert{minimizacija logičkih funkcija}: realizacija mreže sa što manjim brojem komponenti.
\par\vspace{3mm}
% element: 02-s032:e04
To ne mora biti i ekonomski najisplativije. Razvijeni su mnogi algoritmi; mnogi minimizuju funkcije realizovane I i ILI kolima i ukupan broj ulaza u logička kola.
\par\vspace{3mm}
% element: 02-s032:e01
Tehnike minimizacije polaze od funkcije date potpunim proizvodima ili zbirovima. Uočavaju se članovi čiji su \alert{indeksi sa Hemingovim rastojanjem jedan}, odnosno koji se razlikuju samo u jednom literalu.
\par\vspace{3mm}
\Oznaka{Primer}\vspace{2mm}
{\centering \(F=\cdots+C\bar B\bar A+CB\bar A+\cdots\)\par}
\end{frame}
